\section{$\cN{=}2$ 多重态与拉格朗日量}\label{sec:lagrangian}

\subsection{微观拉格朗日量}
\subsubsection{$\cN{=}1$ 超场}\index{Superfield@超场|hyperpage}\index{ChiralMultiplet@手征多重态|hyperpage}\index{AuxiliaryField@辅助场|hyperpage}
现在转向具有 $\cN{=}2$ 超对称性的拉格朗日量的构造。一个 $\cN{=}2$ 超对称理论，特别地也是一个 $\cN{=}1$ 超对称理论。因而，用 $\cN{=}1$ 超场来描述 $\cN{=}2$ 系统很方便。为此，先简要回顾 $\cN{=}1$ 形式体系。仅在本节中，我们将虚数单位写成 $\ii$，以示区别。

一个 $\cN{=}1$ 矢量多重态由一个 Weyl 费米子 $\lambda_\alpha$ 和一个矢量场 $A_\mu$ 组成，二者都属于规范群 $G$ 的伴随表示。将它们合并为超场 $W_\alpha$，其展开为
\begin{equation}
W_\alpha=\lambda_\alpha+F_{(\alpha\beta)}\theta^\beta + D\theta_\alpha+\cdots
\end{equation}
其中 $D$ 是辅助场，同样属于规范群的伴随表示。$F_{\alpha\beta}=\frac \ii2 \sigma^\mu{}^\beta_{\dot\gamma} \bar\sigma^{\nu}{}^{\dot \gamma}_\alpha F_{\mu\nu}$ 是场强 $F_{\mu\nu}$ 的反自对偶部分。

矢量多重态的动能项为
\begin{equation}
\int d^2\theta \frac{-\ii}{8\pi} \tau \tr W_\alpha W^\alpha + cc.
\end{equation}
其中
\begin{equation}
\tau=\frac{4\pi \ii}{g^2}+\frac{\theta}{2\pi}
\end{equation}
是将耦合常数的倒数与 $\theta$ 角合并而成的复数。我们将其称为规范多重态的复化耦合。按分量展开，得到
\begin{equation}
\frac{1}{2g^2} \tr F_{\mu\nu}F^{\mu\nu} + \frac{\theta}{16\pi^2} \tr F_{\mu\nu} \tilde F^{\mu\nu} + \frac{1}{g^2}\tr D^2 -\frac{2\ii}{g^2}\tr \bar\lambda\slash\partial \lambda.
\end{equation}
对于规范代数的标准生成元，我们采用约定 $\tr T^a T^b = \frac12 \delta^{ab}$，这就解释了规范动能项前为何出现因子 $1/(2g^2)$。$\theta$ 项是一个规范依赖项的全导数，因此不影响微扰计算。它确实会影响非微扰计算，我们将在后面回到这一点。

一个 $\cN{=}1$ 手征多重态 $Q$ 由一个复标量 $Q$ 和一个 Weyl 费米子 $\psi_\alpha$ 组成，二者属于规范群的同一表示。用超场表示，有
\begin{equation}
Q=Q\big|_{\theta=0} + 2\psi_\alpha \theta^\alpha + F \theta_\alpha \theta^\alpha
\end{equation}
其中 $F$ 为辅助场。中间分量前的系数 2 不合通常惯例，但这一选择使我们能够消去后面公式中出现的各种烦人的 $\sqrt{2}$ 因子。手征多重态 $Q_{1,\ldots}$ 可以属于规范群 $G$ 的任意复表示 $R$。此时，拉格朗日密度为
\begin{equation}
\int d^4\theta Q^\dagger{}^{j} e^{V^a \rho_a{}^{i}_{j}} Q_i + \int d^2\theta W(Q)+ cc.
\end{equation}
其中 $V$ 是矢量超场，$\rho_a{}^i_j$ 是规范代数的矩阵表示，而 $W(Q)$ 是 $Q_{1,\ldots}$ 的规范不变全纯函数。

要求各个场的超对称变换为零，即可得到超对称真空。非平凡条件来自
\begin{equation}
\delta \lambda_\alpha =0, \qquad \delta\psi_\alpha=0
\end{equation}
它们给出
\begin{equation}
D_a=0, \quad F_i=0.
\end{equation}
解出辅助场的代数运动方程，得到
\begin{equation}
Q^\dagger_{\bar j} \rho_a^{\bar ji}Q_i =0,\quad
\frac{\partial W}{\partial Q_i}=0.
\end{equation}

\subsubsection{矢量多重态与超多重态}\label{sec:hyperdef}\index{VectorMultiplet@矢量多重态|hyperpage}\index{Hypermultiplet@超多重态|hyperpage}\index{RSymmetry@R 对称性|hyperpage}\index{HalfHypermultiplet@半超多重态|hyperpage}\index{NFourTheory@$\cN{=}4$ 理论|hyperpage}\index{NTwoStarTheory@$\cN{=}2^*$ 理论|hyperpage}
一个 $\cN{=}2$ 矢量多重态由如下两个 $\cN{=}1$ 多重态组成，二者都属于规范群 $G$ 的伴随表示：
\begin{equation}
\begin{array}{ccccc@{\qquad}c}
&\raisebox{-.8ex}{\rotatebox{30}{$\leftrightarrow$}}& \lambda_\alpha & \leftrightarrow & A_\mu & \text{$\cN{=}1$ 矢量多重态},\\
\Phi & \leftrightarrow & \tilde\lambda _\alpha &  \rotatebox{30}{$\leftrightarrow$} &  & \text{$\cN{=}1$ 手征多重态}.
\end{array}
\end{equation}
这里，水平箭头表示在 $\cN{=}1$ 超场形式体系中显式呈现的 $\cN{=}1$ 子超对称生成元，斜箭头则对应第二组 $\cN{=}1$ 子超对称性。

构造第二组超对称作用的一种简单方法，是要求理论在作用于 $\lambda_\alpha$ 和 $\tilde\lambda_\alpha$ 的 $\SU(2)$ 转动下对称。一般而言，不与超对称生成元对易的对称性称为 R 对称性。因而，这一 $\SU(2)$ 对称性经常被称为 $\SU(2)_R$ 对称性。将超场形式体系中显式的超对称性与某个 R 对称性结合起来，构造具有更多超对称性的理论，如今已经是一种标准技巧；有关应用，可参见文献\cite{Aharony:2008ug}。还应当记住，不具有 $\SU(2)_R$ 对称性的 $\cN{=}2$ 超对称理论是可能存在的，而且确实存在：可以只有两组超对称生成元，而没有联系它们的 $\SU(2)$ 对称性，例如参见文献\cite{Antoniadis:1995vb,FujiwaraItoyamaSakaguchi}。尽管如此，为简单起见，本讲义只讨论具有 $\SU(2)_R$ 对称性的 $\cN{=}2$ 超对称系统。

此时拉格朗日量为
\begin{equation}
\frac{\Im\tau}{4\pi} \int d^4\theta \tr \Phi^\dagger e^{[V,\cdot]} \Phi +\int d^2\theta \frac{-\ii}{8\pi}\tau \tr W_\alpha W^\alpha + cc.\label{vectorlag}
\end{equation}
凯勒势与规范动能项前的系数之比，由 $\SU(2)_R$ 对称性的要求确定。

一个 $\cN{=}2$ 超多重态\footnote{有一个愚蠢的惯例：拼写“vector multiplets”时，要在“vector”和“multiplets”之间加空格，而拼写“hypermultiplets”时却不加空格。在口语中，hypermultiplets 经常简称为 hypers。}包含如下各场：
\begin{equation}
\begin{array}{ccccc@{\qquad}c}
&\raisebox{-.8ex}{\rotatebox{30}{$\leftrightarrow$}}& Q & \leftrightarrow & \psi & \text{$\cN{=}1$ 手征多重态}\\
\tilde\psi^\dagger & \leftrightarrow &\tilde Q^\dagger &   \rotatebox{30}{$\leftrightarrow$}  &  & \text{$\cN{=}1$ 反手征多重态}
\end{array}
\end{equation}
它们都属于规范群的同一表示 $R$。因而，$\cN{=}1$ 手征多重态 $Q$ 和 $\tilde Q$ 属于规范群的共轭表示。为使理论具有 $\cN{=}2$ 超对称性，我们再次要求它在作用于 $Q$ 和 $\tilde Q^\dagger$ 的 $\SU(2)$ 转动下对称。

为具体起见，考虑 $G=\SU(N)$ 以及基本 $N$ 维表示中的 $N_f$ 个超多重态 $Q_i^a$、$\tilde Q^i_a$，其中 $a=1,\ldots,N$，$i=1,\ldots,N_f$。这组场经常被称为 $\SU(N)$ 的 $N_f$ 个基本表示味。规范变换对它们的作用为
\begin{equation}
Q_i \to e^{\Lambda} Q_i,\qquad
\tilde Q^i \to  \tilde Q^i e^{-\Lambda}
\end{equation}
其中 $\Lambda$ 是由手征超场组成的无迹 $N\times N$ 矩阵；这里省略了规范指标。

超多重态的拉格朗日量为
\begin{equation}
c\int d^4\theta (Q^\dagger{}^i e^V Q_i+ \tilde Q^i e^{-V} \tilde Q^\dagger{}_i)
+c'(\int d^2 \theta \tilde Q^i \Phi Q_i + cc.)
+(\int d^2\theta \mu^i_j \tilde Q^jQ_i + cc.)\label{hyperlag}
\end{equation}
其中再次省略了规范指标 $a$。$\SU(2)_R$ 对称性的存在确定了 $c$ 与 $c'$ 的比值：例如，可以比较第一项中 $Q^i\lambda\psi$ 的系数和第二项中 $\tilde Q^i\tilde\lambda\psi$ 的系数。我们发现，取 $c=c'$ 即可。除非另有说明，以下均取 $c=c'=1$。$\SU(2)_R$ 对称性还要求质量项 $\mu^i_j$ 满足 $[\mu,\mu^\dagger]=0$。因此，$\mu$ 可以对角化，质量项也就经常写为
\begin{equation}
\sum_i \int d^2\theta \mu_i \tilde Q^i Q_i + cc.
\end{equation}

作为另一个例子，考虑伴随表示中的一个超多重态 $(Z,\tilde Z)$，即二者都是 $N\times N$ 无迹矩阵。以下讨论也很容易推广到任意规范群。当超多重态无质量时，总拉格朗日量具有如下形式：
\begin{multline}
\int d^2\theta \frac{-\ii}{8\pi}\tau \tr W_\alpha W^\alpha + cc.
+\frac{\Im\tau}{4\pi} \int d^4\theta \tr \Phi^\dagger e^{[V,\cdot]} \Phi \\
+\frac{\Im\tau}{4\pi} \int d^4\theta (Z^\dagger{} e^{[V,\cdot]} Z+ \tilde Z e^{-[V,\cdot]} \tilde Z^\dagger{})
+\frac{\Im\tau}{4\pi} \int d^2 \theta \tilde Z [\Phi, Z] + cc.\label{n4lag}
\end{multline}
这里，我们对式\eqref{hyperlag}中的 $c=c'$ 作了不同选择。这一拉格朗日量显然具有转动 $\Phi$、$Z$ 和 $\tilde Z$ 的 $\SU(3)_F$ 味对称性。它与超场形式体系中显式的 $\cN{=}1$ 超对称性对易。我们还知道，该理论具有转动 $Z$ 和 $\tilde Z^\dagger$ 的 $\SU(2)_R$ 对称性。$\SU(3)_F$ 与 $\SU(2)_R$ 这两个对称性不对易：我们发现，存在作用于
\begin{equation}
\Re \Phi, \Im \Phi, \Re Z, \Im Z, \Re \tilde Z, \Im \tilde Z.
\end{equation}
的 $\SO(6)_R$ 对称性。注意，$\SO(6)_R$ 也可以视为 $\SU(4)_R$，因为 $\SO(6)$ 与 $\SU(4)$ 具有相同的李代数。于是，$\SU(4)_R$ 对称性作用于系统中的四个 Weyl 费米子
\begin{equation}
\lambda, \tilde\lambda, \psi, \tilde \psi
\end{equation}
其中 $\lambda$ 和 $\tilde\lambda$ 属于 $\cN{=}2$ 矢量多重态，而 $\psi$、$\tilde\psi$ 属于 $\cN{=}2$ 超多重态。我们得出结论，该系统实际上具有 $\cN{=}4$ 超对称性，其四个超对称生成元受 $\SU(4)_R\simeq\SO(6)_R$ 作用。这里的论证，是结合超场形式体系中显式与非显式对称性的又一次应用。

可以向式\eqref{n4lag}加入质量项 $\int d^2\theta \mu Z\tilde Z + cc.$。它保持 $\cN{=}2$ 超对称性，但破坏 $\cN{=}4$ 超对称性。得到的理论有时称为 $\cN{=}2^*$ 理论。

在结束这一小节之前，还应当提及半超多重态的概念。从一个完整超多重态 $(Q_a,\tilde Q^a)$ 出发，其中 $Q_a$ 与 $\tilde Q^a$ 分别属于表示 $R$ 和 $\bar R$。当 $R$ 为赝实表示，或等价地，当存在反对称不变张量 $\epsilon_{ab}$ 时，可以施加与 $\cN{=}2$ 超对称性相容的约束
\begin{equation}
Q_a=\epsilon_{ab} \tilde Q^b
\end{equation}
使多重态中的自由度数目减半。所得多重态称为表示 $R$ 中的半超多重态。第\ref{sec:hyperrevisited}小节还将回到这一概念。

\subsection{真空}\label{sec:classicalvac}\index{SupersymmetricVacuum@超对称真空|hyperpage}\index{CoulombBranch@库仑支|hyperpage}\index{HiggsBranch@希格斯支|hyperpage}\index{MixedBranch@混合支|hyperpage}\index{DTerm@D 项|hyperpage}\index{FTerm@F 项|hyperpage}

矢量多重态与超多重态组成的系统，其拉格朗日量是式\eqref{vectorlag}与式\eqref{hyperlag}之和。超对称真空由以下条件给出。

首先，对辅助场 $D$ 变分得到
\begin{equation}
\frac{1}{g^2}[\Phi^\dagger,\Phi]+ (Q_i Q^\dagger{}^i-\tilde Q^\dagger{}_i \tilde Q^i)\big|_\text{traceless}=0,\label{Deq}
\end{equation}
其中，对于一个 $N\times N$ 矩阵，$X|_\text{traceless}$ 定义为
\begin{equation}
X|_\text{traceless}=X-\frac1N\tr X.
\end{equation}
我们约定，在需要时将标量乘以单位矩阵。

其次，对 $\Phi$ 的辅助场 $F$ 变分得到
\begin{equation}
Q_i \tilde Q^i \big|_\text{traceless}=0\label{Feq}
\end{equation}
而 $Q_i$、$\tilde Q^i$ 的辅助场 $F$ 给出
\begin{equation}
\Phi Q_i +\mu^j_i Q_j =0, \qquad \tilde Q^i \Phi + \mu^i_j \tilde Q^j=0 \label{Qeq}
\end{equation}
它们对所有 $i$ 都成立。总标量势是式\eqref{Deq}、\eqref{Feq}和\eqref{Qeq}中各表达式的绝对值平方的加权和。

到目前为止，我们只使用了相对于超场记号中显式的 $\cN{=}1$ 超对称性的超对称条件。对式\eqref{Deq}第一项与第二项之间的交叉项作适当整理，再与式\eqref{Qeq}的平方项合并，就可以把总标量势重写为下列各对象的加权和。首先，有一个只含 $\Phi$ 的项：
\begin{equation}
[\Phi^\dagger,\Phi]=0.\label{phieq}
\end{equation}
其次，有只含 $Q$ 和 $\tilde Q$ 的项，其中之一为
\begin{equation}
(Q_i Q^\dagger{}^i-\tilde Q^\dagger{}_i \tilde Q^i)\big|_\text{traceless}=0\label{tripleteq}
\end{equation}
另一个则是式\eqref{Feq}。最后，还有混合 $\Phi$ 和 $Q$ 的项，即式\eqref{Qeq}以及
\begin{equation}
\Phi^\dagger Q_i +\mu^\dagger{}^j_i Q_j =0, \qquad \tilde Q^i \Phi^\dagger + \mu^\dagger{}^i_j \tilde Q^j=0 \label{QQeq}.
\end{equation}

注意，式\eqref{phieq}与式\eqref{tripleteq}分别是方程\eqref{Deq}的 $\SU(2)_R$ 单态部分和三重态部分。此外，方程\eqref{tripleteq}与方程\eqref{Feq}的实部和虚部共同构成 $\SU(2)_R$ 的三重态。最后，方程\eqref{Qeq}与方程\eqref{QQeq}按照 $\SU(2)_R$ 的二重态变换。

总结一下。我们最初在式\eqref{Deq}、\eqref{Feq}和\eqref{Qeq}中要求一组 $\cN{=}1$ 子超对称性不破缺。随后发现，所满足的方程自动具有 $\SU(2)_R$ 不变性，因此全部 $\cN{=}2$ 超对称性也就自动保持不破缺。

获得超对称性的一种简单方式，是要求式\eqref{phieq}成立并令 $Q=\tilde Q=0$。超对称真空模空间的这一子空间称为库仑支，因为通常在红外仍有若干阿贝尔规范场保留下来。

当质量项 $\mu^i_j$ 非零时，一般而言，其他真空组态的讨论并不直接。当 $\mu^i_j=0$ 时，还存在另一类真空组态，只需要求式\eqref{tripleteq}与式\eqref{Feq}成立，并令 $\Phi=0$。这称为希格斯支。这一领域中的一些人只把规范群完全破缺的分支称为希格斯支；但从理论上说，这里所定义的希格斯支在各种操作下具有更统一的行为。

超多重态中的标量 $Q$、$\tilde Q$ 与矢量多重态中的标量 $\Phi$ 同时非零的分支，称为混合支。

由式\eqref{phieq}可知，在超对称真空中，$\Phi$ 可以对角化。为具体起见，令 $G=\SU(2)$。此时 $\Phi=\diag(a,-a)$。当 $a\neq0$ 时，规范群破缺为 $\U(1)$。由于红外仍有库仑场存在，这些真空称为库仑支。下面计算由此产生的 W 玻色子的质量。由
\begin{equation}
\frac1{g^2}\tr|D_\mu \Phi|^2 = \frac1{g^2}\tr (\partial_\mu \Phi+[A_\mu,\Phi])^2
\end{equation}
可知，拉格朗日量中存在一项
\begin{equation}
\frac1{g^2}\tr [A_\mu,\vev{\Phi}]^2
\end{equation}
它赋予矢量场以质量。写成
\begin{equation}
A_\mu=\begin{pmatrix}A^0 & W^+ \\W^- & -A^0\end{pmatrix}_\mu,
\end{equation}
可得
\begin{equation}
\left[
\begin{pmatrix}0 & W^+_\mu \\0 & 0\end{pmatrix},
\begin{pmatrix}a & 0\\0 & -a\end{pmatrix}
\right]=-2a \begin{pmatrix}0 & W^+_\mu \\0 & 0\end{pmatrix}.
\end{equation}
在我们的约定中，动能项为 $\tr F_{\mu\nu}F_{\mu\nu}/(2g^2)$，因此得到质量
\begin{equation}
M_W=|2a|.
\end{equation}
对于固定的 $i$，场 $Q_i$、$\tilde Q^i$ 的质量项为
\begin{equation}
\tilde Q^i \Phi Q_i+ \mu_i \tilde Q^i Q_i =
(\tilde Q_1^i, \tilde Q_2^i)
\begin{pmatrix}a+\mu_i & 0 \\0 & -a+\mu_i\end{pmatrix}
\begin{pmatrix}Q_1^i \\Q_2^i\end{pmatrix}.
\end{equation}
因此有
\begin{equation}
M_{Q_{i,1}}=|a+\mu|,\qquad
M_{Q_{i,2}}=|{-a}+\mu|.
\end{equation}

在 $\theta=0$ 的情形下，我们已经在式\eqref{monopole-energy}中研究了这一模型的单极子的经典质量。一般而言，它由
\begin{equation}
M_\text{monopole}=|2\tau a|.
\end{equation}
给出。经典地，粒子的质量满足一个一般不等式：
\begin{equation}
M \geq |na+m(2\tau a) + \sum_i f_i \mu_i| \label{classicalBPS}
\end{equation}
其中 $n$、$m$、$f_i$ 分别为粒子的电荷、磁荷和味荷。这里，第 $i$ 个味荷对应于对称性
\begin{equation}
Q_i \to e^{\ii\varphi_i} Q_i,\quad
\tilde Q^i \to e^{-\ii\varphi_i} \tilde Q^i.
\end{equation}
这一不等式称为 Bogomolnyi-Prasad-Sommerfield（BPS）界；只要计入 $a$ 与 $2\tau a$ 的量子修正，它在量子系统中仍然成立。接下来研究这一点。

\subsection{BPS 界}\index{BPSBound@BPS 界|hyperpage}\index{BPSState@BPS 态|hyperpage}\index{CentralCharge@中心荷|hyperpage}\index{BPSMassFormula@BPS 质量公式|hyperpage}
一般的 $\cN{=}2$ 超对称代数具有如下形式：
\begin{align}
\{Q^I_\alpha,Q^\dagger{}^{\bar J}_{\dot \beta}\} &= \delta^{I\bar J} P_{\mu} \sigma^\mu_{\alpha\dot\beta},\\
\{Q^I_\alpha , Q^J_\beta\}&= \epsilon^{IJ} \epsilon_{\alpha\beta}Z.\label{Zcomm}
\end{align}
这里，$I=1,2$ 是区分两个超对称生成元的指标，而 $Z$ 是一个与所有量都对易的复量。取满足
\begin{equation}
P_\mu = (M,0,0,0).
\end{equation}
的坐标系。这一选择将洛伦兹对称性 $\SO(3,1)$ 破缺为三维空间转动 $\SO(3)$，使我们能够将不带点和带点的旋量指标等同起来。随后定义
\begin{equation}
{}^{(\varphi)}Q_\alpha=\frac{1}{\sqrt2}(Q^1_\alpha + e^{-\ii\varphi} \sigma^0{}_{\alpha}{}^{\dot\beta} Q^\dagger{}^2_{\dot\beta})
\end{equation}
它满足
\begin{equation}
\{
{}^{(\varphi)}Q_\alpha,
{}^{(\varphi)}Q^\dagger_\beta
\}=\delta_{\alpha\beta} (M-\Re (e^{-\ii\varphi} Z)).\label{varphi}
\end{equation}

一般而言，如果存在一个算符 $a$，满足 $\{a,a^\dagger\}=c$，其中 $c$ 为常数，那么 $c$ 必须非负。事实上，取一个右矢 $\ket{\psi}$，便有
\begin{equation}
\left|a^\dagger \ket{\psi}\right|^2
+\left|a \ket{\psi}\right|^2
=\vev{\psi|a a^\dagger|\psi}+\vev{\psi|a^\dagger a|\psi} = c\vev{\psi|\psi},\label{tmp}
\end{equation}
这意味着 $c\ge0$。于是，由式\eqref{varphi}可知
\begin{equation}
M\ge \Re (e^{-\ii\varphi} Z)
\end{equation}
对所有 $\varphi$ 都成立。选择 $\varphi=\Arg Z$，就得到不等式
\begin{equation}
M\geq |Z|. \label{BPS}
\end{equation}

一般而言，超平移 $Q^I_\alpha$ 和 $Q^J{}^\dagger_\alpha$ 的多重态在超多重态中生成 $2^4=16$ 个态。当不等式\eqref{BPS}取等号时，对于 $a_\alpha={}^{(\Arg Z)}Q_\alpha$，式\eqref{tmp}中的 $c$ 为零，迫使算符 ${}^{(\Arg Z)}Q_\alpha$ 本身为零。于是，超平移只生成 $2^2=4$ 个态。这样的多重态称为 BPS 多重态，而那些在超平移作用下包含 16 个态的多重态称为非 BPS 多重态。BPS 态相当稳健：在一般扰动下，一个多重态中的态数不能发生跃变。因此，一般情况下，BPS 态会继续保持为 BPS 态。

这个与所有量都对易的量 $Z$ 究竟是什么？按照定义，与所有量都对易的量就是守恒荷。当低能理论是弱耦合的 $\U(1)$ 规范理论时，$Z$ 是电荷 $n$、磁荷 $m$ 与味荷 $f_i$ 的线性组合。在量子理论中，我们将这一线性组合中出现的系数定义为 $a$、$a_D$ 和 $\mu_i$：
\begin{equation}
Z=na+ma_D+\sum_i \mu_i f_i.\label{Zlin}
\end{equation}
当理论弱耦合时，通过比较量子 BPS 质量公式\eqref{BPS}与其经典对应式\eqref{classicalBPS}，可以将 $a$ 认定为场 $\Phi$ 的对角矩阵元，将 $a_D$ 认定为 $2\tau a$，并将 $\mu_i$ 认定为拉格朗日量中质量项的系数。在强耦合区域，说 $a$ 是规范依赖场 $\Phi$ 的对角矩阵元并没有意义。相反，应当把式\eqref{Zlin}视为量 $a$ 的定义。

\subsection{低能拉格朗日量}\label{sec:lel}\index{Prepotential@预势|hyperpage}\index{SpecialKahlerManifold@特殊凯勒流形|hyperpage}\index{SpecialCoordinates@特殊坐标|hyperpage}\index{LowEnergyLagrangian@低能拉格朗日量|hyperpage}
考虑一个在红外描述 $\U(1)^n$ 规范场的一般有效拉格朗日量。将 $n$ 个 $\U(1)$ 矢量多重态记为
\begin{equation}
\begin{array}{ccccc@{\qquad}c}
&\raisebox{-.8ex}{\rotatebox{30}{$\leftrightarrow$}}& \lambda_\alpha & \leftrightarrow & A_\mu & \text{$\cN{=}1$ 矢量多重态}\\
a & \leftrightarrow & \tilde\lambda _\alpha & \rotatebox{30}{$\leftrightarrow$}  &  & \text{$\cN{=}1$ 手征多重态}
\end{array}\label{u1mult}
\end{equation}
并附加指标 $i=1,\ldots,n$。一般的 $\cN{=}1$ 超对称拉格朗日量为
\begin{equation}
\frac{1}{8\pi}\int d^4\theta K(\bar a_i,a_j)+\int d^2\theta \frac{-\ii}{8\pi}\tau^{ij}(a) W_{\alpha,i} W^\alpha{}_j + cc.\label{u1kin}
\end{equation}
注意，我们允许凯勒势与规范耦合矩阵非平凡地依赖于 $a_i$。

要求存在转动 $\lambda_\alpha$ 和 $\tilde\lambda_\alpha$ 的 $\SU(2)_R$ 对称性，以保证 $\cN{=}2$ 超对称性的存在。$\tilde\lambda_\alpha$ 的动能矩阵为
\begin{equation}
\frac{1}{4\pi}\frac{\partial^2 K}{\partial a_i \partial\bar a_j}
\end{equation}
而 $\lambda$ 的动能矩阵为
\begin{equation}
\frac{\Im \tau^{ij}}{2\pi}=\frac{\tau^{ij}-\bar \tau^{ij}}{4\pi \ii}.
\end{equation}
令二者相等，得到
\begin{equation}
\frac{\tau^{ij}-\bar \tau^{ij}}\ii =\frac{\partial^2 K}{\partial a_i \partial\bar a_j}.
\end{equation}
两边对 $a_k$ 求导，得到
\begin{equation}
\frac{\partial}{\partial a_k} \frac{\tau^{ij}}\ii=\frac{\partial^3 K}{\partial a_k\partial a_i \partial\bar a_j}.
\end{equation}
左边在 $i\leftrightarrow j$ 下对称，右边在 $k\leftrightarrow i$ 下对称。因此，至少在局部，$\tau^{ij}$ 可以积分两次：
\begin{equation}
\tau^{ij}=\frac{\partial^2 F}{\partial a_i\partial a_j}\label{pretau}
\end{equation}
其中 $F(a)$ 为局部全纯函数。定义
\begin{equation}
a_D^i=\frac{\partial F}{\partial a_i},\label{adi}
\end{equation}
则有
\begin{equation}
K=\ii(\bar a_D^i a_i-\bar a_i a_D^i ).\label{Keq}
\end{equation}

具有这一附加结构的凯勒流形经常称为特殊凯勒流形。在超引力中，会出现稍有不同的结构。为与之区分，这里也称为刚性特殊凯勒流形。同一几何结构还称为 Seiberg-Witten 可积系统，或 Donagi-Witten 可积系统。有关综述，例如可参见文献\cite{Donagi:1997sr,Craps:1997gp,Freed:1997dp}。在这一语境下，场 $a_i$ 与 $a_D^i$ 称为特殊坐标。

记号 $a_i$ 与 $a_D^i$ 的合理性可以如下说明。假设存在超多重态 $Q$、$\tilde Q$，它仅相对于第 $i$ 个矢量多重态带荷。它具有超势
\begin{equation}
W=Qa_i \tilde Q,
\end{equation}
从而给出质量
\begin{equation}
M_Q=|a_i|.
\end{equation}
因此，$a_i$ 确实就是式\eqref{Zlin}中出现的系数。为说明记号 $a_D^i$ 的合理性，先按分量写出玻色部分的拉格朗日量：
\begin{equation}
\frac{\Im \tau^{ij}}{4\pi} \partial_\mu\bar a_i \partial^\mu a_j
+\frac{\Im \tau^{ij}}{8\pi} F_{\mu\nu\ i} F^{\mu\nu}_j
+\frac{\Re \tau^{ij}}{8\pi} F_{\mu\nu\ i} \tilde F^{\mu\nu}_j.
\end{equation}
推广第\ref{ST}小节的论证，对偶电磁场 $F_D$ 为
\begin{equation}
F_D{}_{\mu\nu}^i = \Im \tau^{ij} F_{\mu\nu\ j} + \Re \tau^{ij} \tilde F_{\mu\nu\ j},
\end{equation}
以它表示，规范场的动能项为
\begin{equation}
\frac{1}{8\pi} \left(\Im \tau_D{}_{ij} F_D{}_{\mu\nu}^{i} F_D{}^{\mu\nu\ j}
+\Re \tau_D{}_{ij} F_D{}_{\mu\nu\ i} \tilde F_D{}^{\mu\nu\ j} \right)
\end{equation}
其中
\begin{equation}
\tau_D{}_{ij} = (-\tau^{-1})_{ij}.
\end{equation}
于是得到
\begin{equation}
\frac{1}{4\pi}\Im \tau^{ij} \partial_\mu\bar a_i \partial^\mu a_j
=\frac{1}{4\pi}\Im \tau_D{}_{ij} \partial_\mu\bar a_D^i \partial^\mu a_D^j
\end{equation}
其中 $a_D$ 按式\eqref{adi}定义。这意味着，我们有对偶的 $\cN{=}2$ 多重态
\begin{equation}
\begin{array}{ccccc@{\qquad}c}
&\raisebox{-.8ex}{\rotatebox{30}{$\leftrightarrow$}}& \lambda_D{}_\alpha & \leftrightarrow & A_D{}_\mu & \text{$\cN{=}1$ 矢量多重态}\\
a_D & \leftrightarrow & \tilde\lambda _D{}_\alpha & \rotatebox{30}{$\leftrightarrow$}   &  & \text{$\cN{=}1$ 手征多重态}
\end{array}
\end{equation}
其中 $A_D{}_\mu$ 是上面引入的 $F_D{}_{\mu\nu}$ 的规范势，并另附上标 $i$。

本小节以一种相当间接的方式引入了预势 $F$，即说明 $\U(1)$ 矢量多重态的动能项\eqref{u1kin}应由式\eqref{pretau}与式\eqref{Keq}给出。利用 $\cN{=}2$ 超空间，可以更好地理解这一点，因为已知预势就是 $\cN{=}2$ 超空间中的拉格朗日密度。这与凯勒势给出 $\cN{=}1$ 超空间中的拉格朗日密度的情形类似。

回顾一下，多重态\eqref{u1mult}可以归纳为 $\cN{=}1$ 超场
\begin{equation}
\Phi_i=a_i+2\tilde\lambda _i{}^\alpha \theta_\alpha+\cdots,\qquad
W_i=\lambda_{\alpha\ i} +F_{\alpha\beta}\theta^\beta+\cdots.
\end{equation}
可以引入另一组超坐标 $\tilde\theta_\alpha$，将它们合并起来：
\begin{equation}
\mathbf{\Phi}_i=\Phi_i+2W_{\alpha\ i} \tilde\theta^\alpha
=a_i + 2\tilde\lambda_{\alpha\ i} \theta^\alpha + 2\lambda_{\alpha\ i}\tilde \theta^\alpha +2 F_{\alpha\beta\ i}\theta^{(\alpha} \tilde\theta^{\beta)}+\cdots.
\end{equation}
于是，转动 $\lambda$ 与 $\tilde\lambda$ 的 $\SU(2)$ R 对称性就作用于两组超坐标 $\theta_\alpha$ 和 $\tilde\theta_\alpha$。

现在，取 $n$ 个变量的任意全纯函数 $F(a_1,\ldots,a_n)$，并考虑它在手征 $\cN{=}2$ 超空间上的积分：
\begin{equation}
\int d^2\theta d^2\tilde \theta F(\mathbf{\Phi}_1,\ldots,\mathbf{\Phi}_n) + cc.
\end{equation}
显然，这会给出规范动能矩阵的结构\eqref{pretau}。为得到凯勒势\eqref{Keq}，需要研究 $\cN{=}2$ 超场的约束和辅助场的结构，例如参见文献\cite{D'Hoker:1999ft}的第 2.10 节。非阿贝尔微观作用量\eqref{vectorlag}的预势为 $F(\mathbf{\Phi})=\frac12\tau\tr\mathbf{\Phi}^2$。
